\documentclass[11pt]{article}
\usepackage[T1]{fontenc}
\usepackage{lmodern,textcomp,amsmath,amssymb}
\usepackage[margin=1in]{geometry}
\usepackage[hidelinks]{hyperref}
\setlength{\parindent}{0pt}
\setlength{\parskip}{0.7em}
\pagestyle{empty}
\begin{document}
{\Large\bfseries Manuscript submission}\par
Journal of Graph Theory\hfill 26 October 2026

Dear Editors,

Please consider the manuscript \emph{Every odd antihole is strongly chromatic-choosable} by Deep Bhattacharjee and Ushashi Bhattacharya for publication in \emph{Journal of Graph Theory}.

The manuscript proves that every odd antihole $\overline{C_{2\ell-1}}$, $\ell\ge3$, is strongly $\ell$-chromatic-choosable: its only non-colourable assignments of lists of size $\ell-1$ are the constant assignments. It also gives the resulting extremal characterization in odd-hole-free critical graphs. These are Conjectures 9 and 3 of Fiujlaali and McDonald, whose contribution and deduction are credited. The proof proceeds through complements of odd linear forests, Hall\textquotesingle s theorem, and a lemma on minimal tight sets. The supplied Lean sources formalise the main theorem; the critical-graph deduction is proved in the text. The paper fits the journal\textquotesingle s interest in structural graph theory and list colouring.

An earlier version is publicly available as Hexagon \href{https://hexagonmath.org/2610.00319v1}{2610.00319v1}. The source, verification code and supporting material are archived with the accompanying release:
\begin{center}\url{https://github.com/creelie/odd-antihole/releases/tag/v1.2.0}\end{center}
The submission includes a separate declaration describing AI assistance and the scope of the formal verification.

Thank you for considering the manuscript.

Sincerely,\par
Deep Bhattacharjee\par
Corresponding author\par
\href{mailto:itsdeep@live.com}{itsdeep@live.com}\par
\href{mailto:d.bhattacharjee@erl-forschung.de}{d.bhattacharjee@erl-forschung.de}\par
ORCID: \href{https://orcid.org/0000-0003-0466-750X}{0000-0003-0466-750X}
\end{document}
